% Compiler avec pdfLaTeX :  pdflatex lebenslauf_muster.tex  (2 fois)

\documentclass[11pt,a4paper]{article}

\usepackage[T1]{fontenc}
\usepackage[utf8]{inputenc}
\usepackage[ngerman]{babel}
\usepackage[scaled]{helvet}
\renewcommand{\familydefault}{\sfdefault}
\usepackage{microtype}
\usepackage[a4paper,left=1.8cm,right=1.8cm,top=1.2cm,bottom=1.5cm]{geometry}
\usepackage{xcolor}
\usepackage{tikz}
\usepackage{enumitem}
\usepackage[hidelinks]{hyperref}

\definecolor{shape}{HTML}{4F6FC0}
\definecolor{accent}{HTML}{5B7DB1}
\definecolor{textgray}{HTML}{3A3A3A}
\definecolor{lightgray}{HTML}{8A8A8A}

\pagestyle{empty}
\setlength{\parindent}{0pt}
\color{textgray}

% ---------- Commandes ----------
\newcommand{\sectionhead}[1]{%
  \vspace{1.1em}
  {\large\bfseries\color{accent}\MakeUppercase{#1}}\\[-0.35em]
  {\color{accent!50}\rule{\linewidth}{0.5pt}}\vspace{0.5em}}

% \entry{date}{titre}{organisation}{puces}
\newcommand{\entry}[4]{%
  \noindent
  \begin{minipage}[t]{2.9cm}\raggedright\small\color{lightgray}#1\end{minipage}%
  \hspace{0.3cm}%
  \begin{minipage}[t]{\dimexpr\linewidth-3.2cm\relax}
    {\normalsize\bfseries\color{black!90}#2}\\[-0.1em]
    {\small\color{accent}#3}
    \begin{itemize}[leftmargin=1.1em,itemsep=0.1em,topsep=0.3em,parsep=0pt,label={\tiny$\bullet$}]
      \small #4
    \end{itemize}
  \end{minipage}\par\vspace{0.55em}}

\begin{document}

% ---------- Kopfbereich ----------
\begin{tikzpicture}[remember picture,overlay]
  % Dekoratives blaues Element oben
  \fill[shape!85,rounded corners=8pt]
    ([xshift=7.0cm]current page.north west) --
    ([xshift=13.0cm]current page.north west) --
    ([xshift=11.0cm,yshift=-1.4cm]current page.north west) --
    ([xshift=8.0cm,yshift=-0.9cm]current page.north west) -- cycle;
  % Name
  \node[anchor=north west,inner sep=0pt] at ([xshift=1.8cm,yshift=-2.0cm]current page.north west)
    {\fontsize{24}{28}\selectfont\bfseries\color{black!80}Emmanuel Dakdem Nguena};
  % Kontakt
  \node[anchor=north east,inner sep=0pt,align=right,font=\footnotesize,text=black!65]
    at ([xshift=-1.8cm,yshift=-1.0cm]current page.north east)
    {Ederstraße 6, 35390 Gießen\\
     017643614045\\
     emmanuel.dakdem@yahoo.com\\
     Führerschein B197};
\end{tikzpicture}

\vspace*{2.7cm}

% Profil
{\small\linespread{1.08}\selectfont
Mechatronik- und Maschinenbau-Student mit starkem Fokus auf Entwicklung, Testen und Optimierung
technischer Lösungen. Durch Projekte in den Bereichen Sensorfusion, Robotik und Systemintegration habe ich
Erfahrung in der Konzeption, Validierung und Dokumentation mechatronischer Systeme gesammelt. Meine
Arbeitsweise zeichnet sich durch strukturiertes Vorgehen, analytische Bewertung und fachübergreifende
Zusammenarbeit mit Entwicklung, Fertigung und Qualitätsprozessen aus.\par}

\sectionhead{Berufserfahrung}

\entry{Apr. 2024 --\\ Feb. 2025}
{Werkstudent neue Technologien und Konnektivität -- Aßlar}
{PFEIFFER VACUUM}
{\item Konzeption, Dokumentation und Abstimmung technischer Anforderungen in interdisziplinären Projekten
 \item Unterstützung bei der Einführung von SAP S/4HANA sowie der Anbindung an bestehende Prozesse
 \item Konzeptentwicklung zur kamerabasierten Qualitätssicherung bei einer Laserentgratanlage
 \item Analyse technischer Anforderungen und deren Abgleich mit Prozessen, Systemen und Lieferanten
 \item Strukturierte Dokumentation und Wissensmanagement im Projektkontext}

\entry{Apr. 2023 --\\ Apr. 2024}
{Werkstudent Qualitätskontrolle -- Heuchelheim}
{Schunk Carbon Technology GmbH}
{\item Tätigkeit im Bereich Qualitätssicherung und Prozesskontrolle
 \item Durchführung von Messprotokollen, Dokumentation und Auswertung von Produktionsdaten mit SAP
 \item Analyse technischer Abweichungen und Unterstützung bei Qualitätsprozessen}

\sectionhead{Ausbildung}

\entry{Okt. 2025 --\\ heute}
{Master of Science (Schwerpunkt: Mechatronik \& Robotik) -- Friedberg}
{Technische Hochschule Mittelhessen}
{\item Vertiefung in interdisziplinären Bereichen der Mechatronik und Robotik
 \item Fokus auf die Integration von Mechanik, Elektronik und Informatik; Systemsimulation und KI-Methoden}

\entry{Okt. 2021 --\\ Sep. 2025}
{Bachelor of Engineering (Schwerpunkt: Allgemeiner Maschinenbau) -- Gießen}
{Technische Hochschule Mittelhessen}
{\item Konstruktion und Inbetriebnahme technischer Maschinen und Anlagen
 \item CAD- und Simulationssoftware (Siemens NX, Ansys, SolidWorks), Automatisierungs- und Messtechnik}

\entry{Sep. 2016 --\\ Aug. 2018}
{Allgemeine Hochschulreife (Abitur) -- Yaoundé, Kamerun}
{Gymnasium CSAO}
{\item Naturwissenschaftlicher Schwerpunkt: Mathematik und Physik
 \item Fundierte Kenntnisse in analytischem Denken und technischer Problemlösung}

\sectionhead{Praktika \& Projekte}

\entry{Apr. 2026 --\\ heute}
{Projekt: Klassische und visuelle Odometrieverfahren für ein automatisiertes Modellfahrzeug}
{Technische Hochschule Mittelhessen, Friedberg}
{\item \textbf{Konzeption \& Implementierung:} Modulares Odometriesystem in ROS~2 für das Modellfahrzeug MxCarkit
 \item \textbf{Sensorfusion:} Kopplung von IMU-Daten, Motordrehzahlen und Lenkwinkel zur Zustandsschätzung und Posenberechnung
 \item \textbf{Computer Vision:} Zwei Verfahren zur visuellen Odometrie auf Basis von Stereo-Kameradaten
 \item \textbf{Systemintegration:} Einbindung in die bestehende ROS~2-Architektur über standardisierte Schnittstellen
 \item \textbf{Validierung \& Analyse:} Experimentelle Fahrversuche zur Bewertung von Genauigkeit, Robustheit und Systemgrenzen}

\entry{März 2025 --\\ Mai 2025}
{Praktikant: Analysetool zur Identifikation von Optimierungspotentialen in der Produktion}
{PFEIFFER VACUUM, Aßlar}
{\item Daten- und Prozessanalyse zur Identifikation von Optimierungspotenzialen in der Produktion
 \item Visualisierung von Kennzahlen mit Power~BI, Erstellung von Dashboards und Berichten
 \item Auswertung von Produktionsprozessen und Identifikation von Engpässen
 \item Anwendung von SAP-Transaktionscodes zur Informationsbeschaffung und Prozessabstimmung}

\sectionhead{Kompetenzen}

\noindent\small\textbf{Sprachen:} Französisch (verhandlungssicher), Deutsch (verhandlungssicher), Englisch (verhandlungssicher)\par
\vspace{0.4em}
\noindent\small\textbf{Software:} Siemens NX (CAD, FEM), SolidWorks, Ansys, ROS~2, Python, SAP S/4HANA, MS Office, Codesys, Siemens TIA Portal\par

\end{document}
